%====================================================
%  Disclaimer
%====================================================
\documentclass[
    aps,
    prd,
    twocolumn,
    nofootinbib,
    floatfix,
    superscriptaddress
]{revtex4-2}

\usepackage[T1]{fontenc}
\usepackage[utf8]{inputenc}
\usepackage{lmodern}
\usepackage{amsmath, amssymb, amsfonts}
\usepackage{mathtools}
\usepackage{physics}
\usepackage{bm}
\usepackage{graphicx}
\usepackage{booktabs}
\usepackage{array}
\usepackage{microtype}
\usepackage{enumitem}
\setlist{nosep}
\usepackage[
    colorlinks=true,
    linkcolor=blue,
    citecolor=blue,
    urlcolor=blue
]{hyperref}

%===================================================

\begin{document}

\title{Gamma-ray Analysis for Multi-Messenger Astrophysics: GAMMA}

\author{Joan Alcaide-Núñez}
\affiliation{Ludwig-Maximilian-Universität München (LMU München), Munich, Germany}

\author{Martina Solano Soto}
\affiliation{affil}

\date{\today}

\maketitle

% Abstract
\begin{abstract}
	Abstract.
\end{abstract}

\tableofcontents
\vspace{1em}

%====================================================
\section{Introduction}
The Gamma-ray Analysis for Multi-Messenger Astrophysics (GAMMA) initiative is a project by the CAPIBARA Collaboration with the goal to study the high-energy, gamma-ray phenomena in the universe and launch a CubeSat to observe these events. In this white paper, we provide a set of compelling open problems that motivate the development and launch of such a mission.

In the 2020 decadal survey \cite{missing}, it was pointed out that multi-band, that means acros the electromagnetic spectrum, and multi-messenger, that means combining different "messengers" such as light, GW and neutrino, detections were to be prioritized in the next years. Currently operating flagship gamma-ray observatories have long passed their main operation lifetimes, and new replacements are still in very early phases.

Professionally or student developed CubeSats with miniaturized detector technology have proven to be able to successfully detect gamma-ray bursts (GRBs) \cite{EIRSAT-1, GRBalpha, HERMES, BurstCube}. However, gamma-ray photons cannot be focused as with a traditional optical telescope. A good way around to accurately observing and locating GRBs is to perform triangulation with multiple satellites in orbit. Therefore, having multiple missions with large fields of view and good resolution capabilities is of very much importance for the proper study of this astrophysical phenomena.

Missing clear explanation of what are gamma-ray bursts and why we need to send detectors to orbit?


%====================================================
\section{Multi-Messenger Astronomy}
Multi-messenger astronomy (MMA) is the term that refers to a recently developed appraoch to studying the universe by combining different, independent sources of information. These sources consist of four fundamental 'messengers': electromagnetic radiation, cosmic rays, neutrinos, and gravitational waves.

For centuries, the only 'messenger' used to observe cosmic events was light. Initially, limited to the optical range, observations eventually expanded across the entire electromagnetic spectrum, from radio waves to $\gamma$-rays. Because this spectrum is so vast, it must be divided into smaller bands observed with wavelength-specific telescopes. To gain a deeper understand of astrophysical sources, astronomers routinely combine data from these different bands (multi-band imaging).

Modern astrophysics has increasingle embraced the full multi-messenger approach moving beyond just photons \cite{2020_astro_decadal}. The universe provides critical information through other probes that carry unique signatures of extreme astrophysical phenomena. These include
\begin{itemize}
    \item cosmic rays: charged particles travelling at relativistic speeds, originating from our sun and deep-space events \cite{missing}
    \item neutrinos: nearly massless particles that traverse the universe with minimal interaction, offering pristine and early insights into sources like supernovae, active galactic nuclei (AGN) and black holes \cite{missing}
    \item gravitational waves (GWs): ripples in the fabric of space-time generated by the acceleration of massive, super-dense objects (black holes and neutron stars), detectable by highly sensitive interferometers \cite{missing}
\end{itemize}

A paradigm shift in this field occurred on 17 August 2017 (GW170817), when a gravitational wave detection of a neutron star merger was directly followed by the observation of a gamma-ray burst (GRB) \cite{Goldstein2017}. From that moment on, the scientific community's goal has been to synthesize these different types of signals to construct a comprehensive picture of the cosmos. Relying on photons alone is no longer sufficient, observing extreme energy conditions requires the combined "senses" of all available messengers to get the full perspective.


Since the detection of GW170817, there have not been more GW-light ulti-messenger events confirmed, despite the massive amount of raw data being analyzed by GW observatories like LIGO, VIRGO and KAGRA. In 2017, neutrino and photonic observations were successfully combined to identify blazars as high-energy sources \cite{IceCube2018Multimessenger}. The scientific community expects to drastically accelerate these discoveries in the coming nyears with next-generation observatories like Einstein Telescope, LISA, IceCube-Gen2.


\subsection{The GAMMA Mission}
GAMMA is designed to contribute directly to this shared goal through $\gamma$-ray observations. Specifically, GAMMA addresses the gap in high-energy transient monitoring. Legacy instruments, such as the Fermi Gamma-ray Space Telescope (launched in 2009), have long surpassed their original operational lifetimes. As these aging flagships eventually retire, there will be critical need for continuous, agile coverage of the $\gamma$-ray sky.

$\gamma$-rays represent the most energetic band of the electromagnetic spectrum and the immediate signature of extreme events, exactly as demonstrated by GW170817. Rapid follow-up observations in this specific band are therefore highly impactful for identifying electromagnetic counterparts to gravitational wave triggers.

Ultimately, multi-messenger astronomy is opening new frontiers, but it requires constant evolution of instrumentation. A focused, student-led mission like GAMMA offers the flexibility and agility needed o provide targeted, high-value data, maintaining the rigorous scientific standards required to participate in modern astrophysical research.
%====================================================
\section{Cosmology: GRBs as High-$z$ Probes}


\begin{itemize}
	\item intro to cosmology: the LambdaCDM model as the model of the universe (explaining accelerated expansion from EFE)
	\item explain the problem of the Hubble tension (what is wrong with our current standard model of the universe?)
	\item how GRB can be used to probe the high-z redshift, why does it matter to probe high-z and what do we need in terms of instrumentation
\end{itemize}


%====================================================
\section{Lorentz Invariance and High-Energy Photon Lags}

\begin{itemize}
	\item briefly explain why a theory of quantum gravity is of such relevance
	\item elaborate on how detecting high-energy photon lags provides direct proof of a granular space-time
	\item transform theory into timing requirements
\end{itemize}


%====================================================
\section{Summary of technical requirements}

% insert table


%====================================================
\section{Conclusions}

% Acknowledgments

\begin{acknowledgments}
	We are grateful to the members of the CAPIBARA Collaboration for useful discussions and feedback. The authors are happy to receive feedback: \href{https://capibara3.github.io/contact.html}{contact us}.
\end{acknowledgments}

% Bibliography
\bibliographystyle{apsrev4-2}
\bibliography{references}

\end{document}
